Large language models (LLMs) such as GPT-4, Claude 3, Llama 3, Gemini 1.5, Qwen, and DeepSeek-LLM now power consumer assistants, clinical decision support, software engineering pipelines, scientific discovery, and autonomous agents. The same generality that makes them useful---open-ended instruction following, retrieval and synthesis, and tool use---also exposes a vast attack surface. \emph{LLM safety} denotes the property that a system declines content liable to cause physical, psychological, financial, societal, or systemic harm while remaining helpful on benign requests. This survey organizes the field along seven dimensions: (i) scope---training-, inference-, and deployment-time risks; (ii) taxonomy---jailbreaks, indirect prompt injection, backdoors, poisoning, hallucination, bias, privacy leakage, dual-use uplift; (iii) historical arc---from PPO (2017) and summarization RLHF (2020) through InstructGPT (2022) and Constitutional AI (2022) to GCG (2023) and Constitutional Classifiers (2025); (iv) key methods---RLHF-PPO, DPO, Safe RLHF, RLAIF, SmoothLLM, Llama Guard, StruQ/SecAlign, RMU unlearning; (v) benchmarks---HH-RLHF, BeaverTails, AdvBench, HarmBench, TrustLLM, DecodingTrust, WMDP, XSTest, AgentDojo; (vi) limitations---shallow alignment, cost asymmetry, multilingual and multimodal gaps, weak-to-strong oversight; and (vii) predictions---eight falsifiable forecasts for 2026--2028 in Section 12.3. Across nine years, more than 600 jailbreak papers, 300+ open safety datasets cataloged by Röttge\ldots{}
